% !PRINTED 155
\weberpage{155}

%<<CONT
$a_0 = 1, \quad a_1 = 2\,a, \quad a_2 = a^2 - 4\,c, \quad a_3 = -\,b^2$

then it follows that
\weq{12}{
\begin{aligned}
D ={}& -\,4\,a^3 b^2 + 144\,a c b^2 + 16\,a^4 c \\
& - 128\,a^2 c^2 + 256\,c^3 - 27\,b^4.
\end{aligned}
}

If, on the other hand, one applies formula (8) to form the discriminant of the cubic equation (11), one obtains

\weq{13}{
27\,D = 4\,(a^2 + 12\,c)^3 - (2\,a^3 - 72\,a c + 27\,b^2)^2.
}

From this one derives the discriminant of the general function of the fourth degree

\weq{14}{
f(x) = a_0\,x^4 + a_1\,x^3 + a_2\,x^2 + a_3\,x + a_4
}

by expressing $a$, $b$, $c$, according to \S.\,34, (4), by

\[
\frac{-\,3\,a_1^2}{8} + a_2
\]

\[
\frac{a_1^3}{8} - \frac{a_1 a_2}{2} + a_3
\]

\[
\frac{-\,3\,a_1^4}{256} + \frac{a_1^2 a_2}{16} - \frac{a_1 a_3}{4} + a_4
\]

then replacing $a_1$, $a_2$, $a_3$, $a_4$ by

\[
\frac{a_1}{a_0}, \quad \frac{a_2}{a_0}, \quad \frac{a_3}{a_0}, \quad \frac{a_4}{a_0}
\]

and multiplying by $a_0^6$.

If, after carrying out this simple calculation, one sets for abbreviation

\weq{15}{
A = a_2^2 - 3\,a_1 a_3 + 12\,a_0 a_4,
}

\weq{16}{
B = 27\,a_1^2 a_4 + 27\,a_0 a_3^2 + 2\,a_2^3 - 72\,a_0 a_2 a_4 - 9\,a_1 a_2 a_3,
}

then one obtains from (13)

\weq{17}{
27\,D = 4\,A^3 - B^2.
}

The quantities $A$ and $B$, which will occupy us further in the next section, are called the first and second invariant of the biquadratic function $f$. By (17), then, the discriminant can be expressed rationally through these quantities, though not in such a way that the coefficients become integral; if one carries out the calculation prescribed in (17), in order to represent $D$ through the coefficients $a_0, a_1, a_2, a_3$ themselves, % [sic] printed coefficient list stops at a_3 although (14) contains a_4
the factor 27 must still cancel out. We shall not set down here the long expression, whose computation offers no difficulty whatever,
%CONT>>
